\section{Limitations}

The license map (P2L)~\cite{jahanshahi2024license} used for our study 
relies on detecting license files committed to a
project's repository. While this approach does not assume that
licenses are consistently and accurately recorded in dedicated
license.md files, it does not check the content of all
files. In practice, licenses might be
specified within individual source files, potentially leading to
underreporting or misclassification of a project's licensing status.

Additionally, while the P2L map captures license changes over time,
it may not fully account for licenses that were temporarily removed
or altered before being reinstated, which could introduce
inaccuracies in assessing a project's long-term licensing practices.
The reliance on the latest status of licenses might also obscure
important historical context, such as changes in licensing
throughout a project's lifecycle, which could be crucial for
understanding the project's evolution and the factors influencing
license choices.

Despite the efforts of manual verification employed in P2L map,
these limitations underscore the need for caution when interpreting
results based on it. Complementary methods, such as
cross-referencing with other data sources, may be necessary to
obtain a more complete and accurate picture of licensing practices
in open source software projects.

\paragraph{Regression-cohort scope.} The within-project regression
analysis is conditioned on the 23{,}619 projects that adopted exactly
one license type initially, exactly one license type at the latest
recorded state, and where these differ ($\sim$0.4\% of the
131M-project ecosystem). Conclusions about license transitions
therefore apply to a self-selected slice of projects that both reach
licensing and then change it; they should not be read as describing
the modal OSS project.

\paragraph{Popularity proxies.} We control for project popularity using
two pre-treatment dependency-graph measures (upstream and downstream
project counts). Platform-visibility signals such as fork and star
counts are not included; incorporating them is left to future work.

\paragraph{Multiple comparisons.} We report seven outcomes across two
designs with several language interactions. The two global effects we
emphasize---commits and sustained activity under
restrictive$\rightarrow$permissive transitions---survive a Bonferroni
correction across the seven outcomes ($\alpha/7 \approx 0.007$; commits
$p<0.002$, active months $p<0.001$) in both designs. We nonetheless
interpret the smaller language-specific effects with appropriate caution
given the number of comparisons.
